%% kapa: exam=1 round=35 year=2024 period=2 form=A subject=appraisal_law subject_ko=감정평가관계법규 seq=4 number=4
%% source: Q-Net 감정평가사 기출문제 (cst00309) artlSeq=5213229 file=raw/1008/5213229/2교시 문제지 원본.pdf
%% conversion: kordoc 4.24.51 -> tools/gen1.py (auto); review=proofread-sonnet; flags=none
\begin{question}{4}
국토의 계획 및 이용에 관한 법령상 용도지역ㆍ용도지구의 내용으로 옳지 않은 것은?
\begin{choices}
  \item 제2종일반주거지역: 중고층주택을 중심으로 편리한 주거환경을 조성하기 위하여 필요한 지역
  \item 일반상업지역: 일반적인 상업기능 및 업무기능을 담당하게 하기 위하여 필요한 지역
  \item 생산녹지지역: 주로 농업적 생산을 위하여 개발을 유보할 필요가 있는 지역
  \item 시가지방재지구: 건축물ㆍ인구가 밀집되어 있는 지역으로서 시설 개선 등을 통하여 재해 예방이 필요한 지구
  \item 집단취락지구: 개발제한구역안의 취락을 정비하기 위하여 필요한 지구
\end{choices}
\end{question}
\answer{1}
